\documentclass[b4paper,7pt,twocolumn,landscape]{jsarticle}
\usepackage[utf8]{inputenc}
\usepackage{amsmath, amssymb}
\usepackage{enumitem}
\usepackage[top=10mm,bottom=10mm,left=10mm,right=10mm]{geometry}
\pagestyle{empty}

% 行間とスペースを詰める設定
\setlength{\baselineskip}{7.5pt}
\setlength{\parskip}{0pt}
\setlength{\parindent}{0pt}
\setlength{\columnsep}{8mm}
\setlist[enumerate]{itemsep=0pt, parsep=0pt, topsep=1pt, leftmargin=15pt, label=\arabic*.}

\begin{document}

\twocolumn[
\begin{center}
\large{\textbf{情報リテラシー 2025年度後期中間試験 問題用紙}}
\end{center}



\noindent
出席番号：\underline{\hspace{2cm}} \quad 氏名：\underline{\hspace{6cm}}

\vspace{1.5mm}
]

\subsubsection*{注意事項}
{
\begin{itemize}
  \item マークシートの学年は、\textbf{1}を記入・マークしてください。
  \item マークシートのクラスは、M:01、E:02、I:03、C:04、A:05の該当する数字を記入・マークしてください。
  \item マークシートの番号は、出席番号を\textbf{3桁}に変換して記入・マークしてください。
  \item 問題は全部で\textbf{50問}あります。\textbf{100点満点}（各2点）です。
  \item 解答は\textbf{マークシート}の解答欄に\textbf{正確に濃く記入}してください。
  \item すべて4択問題です。各設問について最も適切な選択肢を1つ選び、マークしてください。
  \item 試験時間は\textbf{50分}です。
\end{itemize}
}

\subsubsection*{[到達目標1]ネットワーク1（問1〜問7）}

\begin{enumerate}
\item LANとはどのようなネットワークを指すか。\\
1.~国や大陸など広い範囲を結ぶ \quad 2.~家庭や学校など限られた範囲を結ぶ \quad 3.~インターネット全体 \quad 4.~無線通信のみ

\item LAN内部の端末同士を接続し、データを転送する装置はどれか。\\
1.~モデム \quad 2.~ルータ \quad 3.~スイッチ \quad 4.~アクセスポイント

\item TCP/IPモデルのうち、Webページなどのデータをやり取りする層はどれか。\\
1.~トランスポート層 \quad 2.~インターネット層 \quad 3.~アプリケーション層 \quad 4.~ネットワークインターフェース層

\item 「ドメイン名とIPアドレスを対応づける」プロトコルはどれか。\\
1.~HTTP \quad 2.~TCP \quad 3.~DNS \quad 4.~IP

\item TCPの特徴として正しいものはどれか。\\
1.~到達確認を行う \quad 2.~速度重視 \quad 3.~音声通話用 \quad 4.~エラー検出しない

\item 「192.168.1.10」はどの種類のIPアドレスか。\\
1.~グローバルIP \quad 2.~プライベートIP \quad 3.~MACアドレス \quad 4.~サブネットマスク

\item サブネットマスクの主な役割はどれか。\\
1.~データの暗号化 \quad 2.~通信速度の調整 \quad 3.~ネットワーク部とホスト部の区別 \quad 4.~IPアドレスの変換
\end{enumerate}



\subsubsection*{[到達目標1]ネットワーク2（問8〜問14）}

\begin{enumerate}
\setcounter{enumi}{7}
\item HTTPSの特徴として正しいものはどれか。\\
1.~通信内容が暗号化されない \quad 2.~通信内容が暗号化される \quad 3.~高速通信のみ対応 \quad 4.~IPアドレスを変換する

\item URLの中で「www.google.com」の部分は何を示すか。\\
1.~プロトコル \quad 2.~ドメイン名 \quad 3.~パス \quad 4.~パラメータ

\item 「?q=network」のようなURLの末尾の部分は何を表すか。\\
1.~パス \quad 2.~ドメイン \quad 3.~パラメータ \quad 4.~サーバー名

\item メール送信に使われるプロトコルはどれか。\\
1.~POP3 \quad 2.~SMTP \quad 3.~IMAP \quad 4.~HTTP

\item 受信メールをサーバーに保存したまま複数の端末で閲覧できる方式はどれか。\\
1.~SMTP \quad 2.~POP3 \quad 3.~IMAP \quad 4.~DNS

\item 100Mbpsの通信速度で800MBの動画を転送するのにかかる時間として最も適切なものはどれか。ただし、実効速度は理論値の約70\%とする。\\
1.~約64秒 \quad 2.~約80秒 \quad 3.~約91秒 \quad 4.~約128秒

\item ZIP圧縮で最も圧縮効果が高いファイル形式はどれか。\\
1.~JPEG \quad 2.~MP3 \quad 3.~テキストファイル \quad 4.~H.264/AVC
\end{enumerate}



\subsubsection*{[到達目標3]サイバーセキュリティ1（問15〜問21）}

\begin{enumerate}
\setcounter{enumi}{14}
\item 情報セキュリティの3要素「CIA」に含まれないものはどれか。\\
1.~機密性 \quad 2.~完全性 \quad 3.~可用性 \quad 4.~公平性

\item 「情報が改ざんされていない、正確である」ことを示すのはどの要素か。\\
1.~機密性 \quad 2.~完全性 \quad 3.~可用性 \quad 4.~信頼性

\item ランサムウェアの主な特徴として最も適切なものはどれか。\\
1.~ファイルを暗号化して金銭を要求する \quad 2.~ネットワークを高速化する \quad 3.~ファイルを圧縮して送信する \quad 4.~メールを自動削除する

\item 多要素認証に含まれない要素はどれか。\\
1.~知識要素（パスワードなど） \quad 2.~所持要素（スマホなど） \quad 3.~生体要素（指紋など） \quad 4.~通信環境要素（Wi-Fi接続）

\item フィッシングメールへの適切な対応として正しいものはどれか。\\
1.~メール内のリンクをすぐ開く \quad 2.~添付ファイルを確認のため開く \quad 3.~公式サイトで直接確認する \quad 4.~送信者に返信して確認する

\item ファイアウォールの役割として最も適切なものはどれか。\\
1.~不正通信を遮断する \quad 2.~データを暗号化する \quad 3.~通信速度を上げる \quad 4.~ウイルスを削除する

\item VLAN（Virtual LAN）の主な目的として最も適切なものはどれか。\\
1.~ネットワークを高速化する \quad 2.~ネットワークを論理的に分離する \quad 3.~暗号化を行う \quad 4.~ファイルを圧縮する
\end{enumerate}



\subsubsection*{[到達目標3]サイバーセキュリティ2（問22〜問28）}

\begin{enumerate}
\setcounter{enumi}{21}
\item 暗号化の主な目的として最も適切なものはどれか。\\
1.~データ量を減らす \quad 2.~通信速度を上げる \quad 3.~情報の機密性を保護する \quad 4.~通信相手を認証する

\item 共通鍵暗号方式の特徴として正しいものはどれか。\\
1.~公開鍵で暗号化し秘密鍵で復号する \quad 2.~同じ鍵で暗号化と復号を行う \quad 3.~鍵を使わずに暗号化する \quad 4.~送信と受信で別の鍵を使う

\item 公開鍵暗号方式の利点として最も適切なものはどれか。\\
1.~処理が高速である \quad 2.~鍵の配送が不要で安全 \quad 3.~鍵を使わずに暗号化できる \quad 4.~送信者と受信者が同じ鍵を使う

\item デジタル署名の役割として正しくないものはどれか。\\
1.~改ざん検出 \quad 2.~本人確認 \quad 3.~否認防止 \quad 4.~データ圧縮

\item SSL/TLSの役割として正しいものはどれか。\\
1.~通信内容の暗号化と認証 \quad 2.~通信速度の向上 \quad 3.~データ圧縮 \quad 4.~ファイル共有

\newpage

\item 電子透かしの主な目的として最も適切なものはどれか。\\
1.~通信エラーを防ぐ \quad 2.~著作権保護や不正コピーの追跡 \quad 3.~通信を暗号化する \quad 4.~データ圧縮を行う

\item 偶数パリティの説明として最も適切なものはどれか。\\
1.~データの1の数を偶数にそろえる方式 \quad 2.~データの0の数を偶数にそろえる方式 \quad 3.~データを2進数に変換する方式 \quad 4.~エラーを自動的に修正する方式
\end{enumerate}



\subsubsection*{[到達目標5]情報の加工2（問29〜問35）}

\begin{enumerate}
\setcounter{enumi}{28}
\item Excelでファイルをローカルに保存する利点として最も適切なものはどれか。\\
1.~同期に時間がかかる \quad 2.~オフラインでも作業できる \quad 3.~他人と自動共有できる \quad 4.~保存場所が不明になる

\item データ取得の際に最も重要な点はどれか。\\
1.~データ件数をできるだけ多くする \quad 2.~データの信頼性と出典を確認する \quad 3.~データ形式をそろえる \quad 4.~無料データのみを使う

\item オープンデータの特徴として正しいものはどれか。\\
1.~研究目的でしか利用できない \quad 2.~政府や自治体などが公開し、自由に利用できるデータである \quad 3.~登録が必要である \quad 4.~再配布は禁止されている

\item 「身長」「体重」「金額」「時間」はどの尺度に当たるか。\\
1.~名義尺度 \quad 2.~順序尺度 \quad 3.~間隔尺度 \quad 4.~比例尺度

\item 平均値を求めるExcel関数として正しいものはどれか。\\
1.~=MEDIAN() \quad 2.~=MODE.SNGL() \quad 3.~=AVERAGE() \quad 4.~=STDEV.S()

\item 標準偏差を求める目的として最も適切なものはどれか。\\
1.~データのばらつきを表すため \quad 2.~平均値を求めるため \quad 3.~データを分類するため \quad 4.~最大値と最小値を求めるため

\item 相関係数が0.9のときの関係として最も適切なものはどれか。\\
1.~強い正の相関がある \quad 2.~強い負の相関がある \quad 3.~相関がない \quad 4.~データの最大値を示す
\end{enumerate}



\subsubsection*{[到達目標5]データベース（問36〜問42）}

\begin{enumerate}
\setcounter{enumi}{35}
\item データベースの主な目的として最も適切なものはどれか。\\
1.~データを一時的に保存する \quad 2.~データを効率的に整理・検索・管理する \quad 3.~データを暗号化して保存する \quad 4.~データの削除を自動化する

\item リレーショナルデータベースの特徴として正しいものはどれか。\\
1.~表形式でデータを管理する \quad 2.~ファイル単位で保存する \quad 3.~非構造データを扱う \quad 4.~AIが自動分類する

\item 主キー（Primary Key）の役割として正しいものはどれか。\\
1.~データを暗号化する \quad 2.~レコードを一意に識別する \quad 3.~データを分類する \quad 4.~値の平均を求める

\item SQLのSELECT文で「すべての列」を取得する正しいSQL文はどれか。\\
1.~\texttt{SELECT ALL FROM products;} \quad 2.~\texttt{SELECT * FROM products;} \quad 3.~\texttt{SELECT name FROM *;} \quad 4.~\texttt{SELECT FROM products ALL;}

\item SQL文で価格が10000円以上の商品を抽出する正しい条件はどれか。\\
1.~\texttt{WHERE price = 10000;} \quad 2.~\texttt{WHERE price < 10000;} \quad 3.~\texttt{WHERE price >= 10000;} \quad 4.~\texttt{WHERE price <= 10000;}

\newpage

\item ORDER BY句の役割として最も適切なものはどれか。\\
1.~条件を指定して抽出する \quad 2.~データを並べ替える \quad 3.~件数を制限する \quad 4.~集計を行う

\item LIMIT句の役割として正しいものはどれか。\\
1.~抽出条件を指定する \quad 2.~件数を制限する \quad 3.~列の並び順を指定する \quad 4.~列名を省略する
\end{enumerate}

\subsubsection*{[到達目標2]アルゴリズム（問43〜問50）}
\begin{enumerate}
\setcounter{enumi}{42}
\item アルゴリズムの説明として最も適切なものはどれか。\\
1.~データを保存する方法のこと \quad 2.~問題を解決するための明確な手順のこと \quad 3.~コンピュータのプログラミング言語のこと \quad 4.~インターネットに接続する技術のこと
\item アルゴリズムの3つの基本構造に含まれないものはどれか。\\
1.~順次 \quad 2.~分岐 \quad 3.~反復 \quad 4.~並列
\item フローチャートの記号について、正しい組み合わせはどれか。\\
1.~ひし形 - 条件分岐 \quad 2.~角丸四角形 - 処理 \quad 3.~長方形 - 判断 \quad 4.~平行四辺形 - 開始・終了
\item 整列(ソート)アルゴリズムの目的として最も適切なものはどれか。\\
1.~データを削除すること \quad 2.~データを一定の順序に並べ替えること \quad 3.~データを圧縮すること \quad 4.~データを暗号化すること
\item バブルソートの説明として最も適切なものはどれか。\\
1.~隣り合う要素を比較して入れ替える \quad 2.~未整列部分から最小値を選択して先頭に移動する \quad 3.~整列済み部分に新しい値を挿入する \quad 4.~データを半分に分割して探索する
\item 選択ソートの説明として最も適切なものはどれか。\\
1.~隣り合う要素を比較して入れ替える \quad 2.~整列済み部分に新しい値を挿入する \quad 3.~未整列部分から最小値を選択して先頭に移動する \quad 4.~データの中央値を基準に分割する
\item 挿入ソートの特徴として最も適切なものはどれか。\\
1.~ほぼ並んでいるデータに強い \quad 2.~常にすべてのアルゴリズムの中で最速である \quad 3.~データ量が多いほど効率が良くなる \quad 4.~整列済みデータには使用できない
\item アルゴリズムを視覚的に確認できるWebツールとして授業で紹介されたものはどれか。\\
1.~VisuAlgo \quad 2.~GitHub \quad 3.~Google Colab \quad 4.~Scratch
\end{enumerate}


\end{document}
%%%%%%%%%%%%%%%%%%%%%%




情報リテラシー 2025年度後期中間試験 正解一覧
============================================================

【第1回：ネットワーク基礎（問1〜問7）】
問1: 2 (家庭や学校など限られた範囲)
問2: 3 (スイッチ)
問3: 3 (アプリケーション層)
問4: 3 (DNS)
問5: 1 (到達確認を行う)
問6: 2 (プライベートIP)
問7: 3 (ネットワーク部とホスト部の区別)

【第2回：インターネット（問8〜問14）】
問8:  2 (通信内容が暗号化される)
問9:  2 (ドメイン名)
問10: 3 (パラメータ)
問11: 2 (SMTP)
問12: 3 (IMAP)
問13: 3 (約91秒)
問14: 3 (テキストファイル)

【第3回：情報セキュリティ基礎（問15〜問21）】
問15: 4 (公平性)
問16: 2 (完全性)
問17: 1 (ファイルを暗号化して金銭を要求する)
問18: 4 (通信環境要素)
問19: 3 (公式サイトで直接確認する)
問20: 1 (不正通信を遮断する)
問21: 2 (ネットワークを論理的に分離する)

【第4回：暗号化技術（問22〜問28）】
問22: 3 (情報の機密性を保護する)
問23: 2 (同じ鍵で暗号化と復号を行う)
問24: 2 (鍵の配送が不要で安全)
問25: 4 (データ圧縮)
問26: 1 (通信内容の暗号化と認証)
問27: 2 (著作権保護や不正コピーの追跡)
問28: 1 (データの1の数を偶数にそろえる方式)

【第5回：データ処理（問29〜問35）】
問29: 2 (オフラインでも作業できる)
問30: 2 (データの信頼性と出典を確認する)
問31: 2 (政府や自治体などが公開し、自由に利用できるデータである)
問32: 4 (比例尺度)
問33: 3 (=AVERAGE())
問34: 1 (データのばらつきを表すため)
問35: 1 (強い正の相関がある)

【第6回：データベース（問36〜問42）】
問36: 2 (データを効率的に整理・検索・管理する)
問37: 1 (表形式でデータを管理する)
問38: 2 (レコードを一意に識別する)
問39: 2 (SELECT * FROM products;)
問40: 3 (WHERE price >= 10000;)
問41: 2 (データを並べ替える)
問42: 2 (件数を制限する)

【第7回：アルゴリズム（問43〜問50）】}\\
問43: 2 (問題を解決するための明確な手順のこと)\\
問44: 4 (並列)\\
問45: 1 (ひし形 - 条件分岐)\\
問46: 2 (データを一定の順序に並べ替えること)\\
問47: 1 (隣り合う要素を比較して入れ替える)\\
問48: 3 (未整列部分から最小値を選択して先頭に移動する)\\
問49: 1 (ほぼ並んでいるデータに強い)\\
問50: 1 (VisuAlgo)

============================================================
【マークシート用正解番号列（問1〜問50）】
2-3-3-3-1-2-3-2-2-3-2-3-3-3-4-2-1-4-3-1-2-3-2-2-4-1-2-1-2-2-2-4-3-1-1-2-1-2-2-3-2-2-2-4-1-2-1-3-1-1

【簡易版（番号のみ）】
問1-10:  2 3 3 3 1 2 3 2 2 3
問11-20: 2 3 3 3 4 2 1 4 3 1
問21-30: 2 3 2 2 4 1 2 1 2 2
問31-40: 2 4 3 1 1 2 1 2 2 3
問41-50: 2 2 2 4 1 2 1 3 1 1
